\newcommand*{\version}{1.0.0}

\documentclass[12pt,
              article,
              oneside
              ]{memoir}

\usepackage{iftex}

\ifluatex
  \usepackage{fontspec}
\fi

\usepackage[
  % latin,
  % english,
  spanish,% <- principal
  es-minimal,
  es-noshorthands,
  es-tabla,
]{babel}

\ifluatex
% LuaLaTeX --------------------
\defaultfontfeatures{Ligatures=TeX, Scale=MatchLowercase}
\usepackage[MainFeatures={medium,
  Ligatures=Common,
  Numbers=OldStyle,
  SmallCapsFeatures = {LetterSpace=5.0}}]
{junicodevf}
%
% \setsansfont{Libertinus Sans}
% (sin las ligaduras de TeX: -{}-help no se vuelve una raya)
\setmonofont[Scale=MatchLowercase, Ligatures=ResetAll]{DejaVu Sans Mono}
% \setmonofont[Scale=MatchLowercase]{Libertinus Mono}
% \setmonofont[Scale=MatchLowercase}]{Fira Code}
\else
% PDFLaTeX ---------------------
\usepackage[T1]{fontenc}
\usepackage[osf]{libertinus}
\usepackage[scaled=0.9]{FiraMono}
% \usepackage[scaled=0.9]{cascadia-code}
% \usepackage[scaled=0.9]{roboto}
\input glyphtounicode
\pdfgentounicode=1
\fi

\usepackage[babel=true, verbose=false, tracking=true, expansion=true, protrusion=true, final, draft=false]{microtype}
\SetTracking{ encoding = *, shape = sc }{ 30 }
\SetTracking[context = notracking, ]{encoding = *}{0}

\usepackage[autostyle=false, style=british]{csquotes}
\usepackage{latexcolors}

\newcommand{\Rojo}{\color[HTML]{8B0000}}
\newcommand*{\rojo}[1]{\textcolor[HTML]{8B0000}{#1}}
\newcommand*{\rojoit}[1]{\textit{\textcolor[HTML]{8B0000}{#1}}}
\newcommand{\Azul}{\color{bluenivaca}}
\newcommand*{\azul}[1]{\textcolor{bluenivaca}{#1}}
\newcommand{\Verde}{\color{darkgreen}}
\newcommand*{\verde}[1]{\textcolor{darkgreen}{#1}}
\newcommand{\Amarillo}{\color{goldenpoppy}}
\newcommand*{\amarillo}[1]{\textcolor{goldenpoppy}{#1}}
\newcommand{\Morado}{\color{byzantium}}
\newcommand*{\morado}[1]{\textcolor{byzantium}{#1}}
\newcommand{\Naranja}{\color{orange}}
\newcommand*{\naranja}[1]{\textcolor{orange}{#1}}

% \usepackage[margin=3cm]{geometry}

\usepackage{graphicx}
\usepackage{booktabs}

\clubpenalty=9999
\widowpenalty=9999
\brokenpenalty=5000
\frenchspacing
\raggedbottom
\midsloppy
\parindent=14pt
\fboxsep=2pt

\renewcommand\thesection{\arabic{section}}
\chapterstyle{reparticle}

\usepackage{enumitem}
\setlist{noitemsep, nosep}
\setlist[description]{font=\itshape}

\newcommand*{\elipsis}{\,.\,.\,.\,}
\newcommand*{\elipsisb}{\textup{[\elipsis]}}

\renewenvironment{quote}%
 {\begin{list}{}%
              {\setlength\rightmargin{0pt}%
              \setlength\leftmargin{\parindent}}%
  \item[]\small\ignorespaces}
 {\unskip\end{list}}


\definecolor{Link}{HTML}{408080}
\usepackage{hyperref}
\hypersetup{
  colorlinks=true,
  linkcolor=blue,
  filecolor=blue,
  citecolor = black,
  urlcolor=Link,
}
\urlstyle{tt}
\usepackage{bookmark}

\setlength{\droptitle}{-2cm}

\title{\bfseries%
  \texttt{py-slowdown} v.\version{} \\
  Manual de usuario}

\author{%
  Nicolás Vaughan\\
  \begin{minipage}[t]{0.8\linewidth}
    \centering\small%
    \href{mailto:nicolasvaughan@proton.me}{\texttt{nicolasvaughan@proton.me}}\\
    \url{https://nicolasvaughan.org}\\
    \url{https://github.com/nivaca/py-slowdown}\\
    \url{https://codeberg.org/nivaca/py-slowdown}\\
  \end{minipage}}


\date{2026-11-24}

% The pictures, made by tools/manual_images.py
\graphicspath{{img/manual/}}
\usepackage{longtable}

% the title and the author in the PDF's own metadata
\hypersetup{pdftitle={py-slowdown: manual de usuario}, pdfauthor={Nicolás Vaughan}}

% the index (makeindex, with the style manual.ist: see latexmkrc)
\makeindex

% menus, keys, and the pictures of the window
\newcommand*{\menu}[1]{\textit{#1}}
\newcommand*{\then}{\,>\,}
\newcommand*{\key}[1]{\texttt{#1}}
\newcommand*{\shot}[3]{%
  \begin{figure}[htbp]
    \centering
    \includegraphics[width=#2\textwidth,height=0.75\textheight,keepaspectratio]{#1}
    \caption{#3}
    \label{fig:#1}
  \end{figure}}

\begin{document}
\maketitle

\tableofcontents*

\section{Introducción}
\noindent
\texttt{py-slowdown} reproduce una grabación más lenta o más rápida sin
cambiarle la altura. Lo hice para la música: para practicar una pieza y
sacarla de oído, donde hay que oír un pasaje una y otra vez a una velocidad
que se pueda seguir. Es un programa para la música, no para transcribir la
voz hablada (clases, entrevistas, dictados); lo que hace además de
reproducir, como la cuenta previa, el tempo o la altura en semitonos, está
pensado para músicos.

\index{forma de onda}\index{velocidad}\index{marcadores}\index{formatos}%
La grabación se ve como una forma de onda, sobre la que avanza la posición
mientras suena. Los mandos de siempre la recorren; la velocidad va de
0,25\,× a 3\,× y se puede cambiar mientras suena; los marcadores, cada uno
con su nombre, guardan los lugares a los que se quiere volver. El programa
recuerda, por archivo, sus marcadores y el punto donde se dejó, y la vez
siguiente sigue desde allí.

\index{rubber band@Rubber Band}\index{calidad}%
El estiramiento lo hace la biblioteca Rubber Band, con su motor R3, el de
más calidad. A 1\,× el archivo suena tal
cual, muestra por muestra, sin pasar por Rubber Band.

\index{idioma de la interfaz}%
La ventana está en inglés o en español. Este manual describe la versión
\version{}; las imágenes son de la ventana en español.

\section{Instalación}
\label{sec:instalacion}

\index{instalacion@instalación}%
Las versiones están en GitHub y en Codeberg (las direcciones están en la
primera página). Cada una trae los instaladores y paquetes que siguen, y
todos los archivos van firmados: véase la sección~\ref{sec:firmas}. Cada
instalador lleva todo lo que el programa necesita, su propio Python y sus
bibliotecas incluidos; no hay que instalar nada más.

\subsection{Windows}

\index{windows@Windows}\index{instalador}%
\texttt{py-slowdown-\version-windows-x64-setup.exe} instala el programa en
Windows 10 o posterior. Lo instala para el usuario que lo corre, sin
permisos de administrador, salvo que en su primera página se elija
instalarlo para todos. Pone una entrada en el menú Inicio, un ícono en el
escritorio si se pide, y \menu{py-slowdown} en \menu{Abrir con} de los
archivos de audio (sin quitarle a ningún otro programa los archivos que
abre); el programa se quita en \menu{Configuración\then Aplicaciones}, como
cualquier otro.

\index{smartscreen@SmartScreen}%
El instalador no está firmado con un certificado que Windows conozca, así
que la primera vez Windows puede decir que protegió el equipo
(SmartScreen). \menu{Más información\then Ejecutar de todas formas} sigue
adelante. La firma de la sección~\ref{sec:firmas} es la que prueba que el
archivo es el publicado.

\subsection{macOS}

\index{macos@macOS}%
Hay una imagen de disco para cada clase de Mac:
\texttt{py-slowdown-\version-macos-arm64.dmg} para los que tienen Apple
Silicon (M1 y posteriores), y \texttt{py-slowdown-\version-macos-x86\_64.dmg}
para los que tienen Intel. Necesitan macOS 13 (Ventura) o posterior. Se
abre la imagen y se arrastra \menu{py-slowdown} a \menu{Aplicaciones}.
Desde entonces los archivos de audio se pueden abrir con él desde el Finder
(\menu{Abrir con}), o soltarlos sobre su ícono.

\index{gatekeeper@Gatekeeper}%
La aplicación no está notarizada por Apple, así que la primera vez macOS
dice que no puede comprobarla. En \menu{Ajustes del Sistema\then
Privacidad y seguridad}, abajo, se elige \menu{Abrir igualmente} para
py-slowdown (en macOS 14 y anteriores, un clic derecho sobre la aplicación
y \menu{Abrir} hace lo mismo). Se pregunta una sola vez.

\subsection{Cualquier Linux: el instalador}

\index{instalador}%
\texttt{py-slowdown-\version-linux-x86\_64.run} instala el programa en
cualquier Linux con glibc 2.34 o posterior, con una entrada en el menú de
aplicaciones:

\begin{quote}
\texttt{sh py-slowdown-\version-linux-x86\_64.run}
\end{quote}

\noindent Lo instala para el usuario que lo corre (en
\texttt{\textasciitilde/.local/share/py-slowdown/app}, con la orden
\texttt{py-slowdown} en \texttt{\textasciitilde/.local/bin}) o, corrido con
\texttt{sudo}, para todos (en \texttt{/opt/py-slowdown}). Pregunta antes
(\texttt{-{}-yes} no pregunta), y \texttt{-{}-check} solo comprueba el
archivo. Para quitar el programa, se corre \texttt{uninstall.sh}, en su
carpeta.

\subsection{Fedora}

\index{fedora (copr, rpm)@Fedora (Copr, RPM)}%
Lo más sencillo, y lo que además lo mantiene al día, es el repositorio de
Copr:

\begin{quote}
\texttt{sudo dnf copr enable nivaca/py-slowdown}\\
\texttt{sudo dnf install py-slowdown}
\end{quote}

\noindent El RPM de cada versión también se puede instalar solo
(\texttt{sudo dnf install ./py-slowdown-\version-1.fc45.noarch.rpm}). Usa
el Rubber Band, el PortAudio y el \texttt{ffmpeg} de Fedora, que
\texttt{dnf} instala con él.

\subsection{Cualquier Linux: la AppImage}

\index{appimage@AppImage}%
La AppImage es un solo archivo que corre en cualquier Linux con glibc 2.34
o posterior (Ubuntu 22.04, Debian 12, Fedora 35 y posteriores) y no
necesita nada instalado: trae su propio Python, Qt, FFmpeg, Rubber Band y
PortAudio. Se vuelve ejecutable y se corre:

\begin{quote}
\texttt{chmod +x py-slowdown-\version-x86\_64.AppImage}\\
\texttt{./py-slowdown-\version-x86\_64.AppImage [ARCHIVO]}
\end{quote}

\subsection{Debian y Ubuntu}

\index{debian y ubuntu@Debian y Ubuntu}%
El paquete de Debian, para Debian 12, Ubuntu 22.04 y posteriores, está
hecho de la AppImage y la instala en \texttt{/opt/py-slowdown}, con la
orden \texttt{py-slowdown}, una entrada en el menú de aplicaciones y la
página de manual:

\begin{quote}
\texttt{sudo apt install ./py-slowdown\_\version-1\_amd64.deb}
\end{quote}

\subsection{Desde el código, en Linux y macOS}

\index{instalador}\index{macos@macOS}%
Con el código (clonado, o el ZIP de una versión desempacado) y Python 3.14
o posterior, \texttt{install/install.sh} instala el programa para un
usuario, en un entorno propio, con una entrada en el menú de aplicaciones
(en macOS, una aplicación en \texttt{\textasciitilde/Applications});
corrido otra vez, lo actualiza, y \texttt{install/install.sh -{}-uninstall}
lo quita. El programa necesita entonces del sistema la biblioteca Rubber
Band (versión 3 o posterior) y, en Linux, PortAudio:

\begin{quote}
Fedora: \texttt{sudo dnf install rubberband-libs portaudio}\\
Debian, Ubuntu: \texttt{sudo apt install librubberband3 libportaudio2}\\
(\texttt{librubberband2} en las versiones anteriores)\\
macOS: \texttt{brew install rubberband}
\end{quote}

\noindent\index{uv}%
Para correrlo desde el código sin instalarlo, con \texttt{uv}:
\texttt{uv run py-slowdown [ARCHIVO]}.

\subsection{Cómo comprobar los archivos de una versión}
\label{sec:firmas}

\index{firmas}%
Cada archivo de una versión tiene su firma (\texttt{ARCHIVO.asc}), y
\texttt{SHA256SUMS} trae las sumas de todos, con su propia firma. Se
comprueban con la clave pública que se publica con la versión:

\begin{quote}
\texttt{gpg -{}-import RPM-GPG-KEY-py-slowdown}\\
\texttt{gpg -{}-verify ARCHIVO.asc ARCHIVO}
\end{quote}

\noindent Los instaladores para Windows y macOS, que hace GitHub, tienen
sus sumas en \texttt{SHA256SUMS-installers}, también firmado. El RPM va
firmado también por dentro; \texttt{dnf} lo comprueba una vez importada la
clave (\texttt{sudo rpm -{}-import RPM-GPG-KEY-py-slowdown}).

\section{Cómo se abre el programa}

\index{arranque}\index{linea de ordenes@línea de órdenes}%
El programa está en el menú de aplicaciones, como \menu{py-slowdown}. Desde
una terminal, \texttt{py-slowdown} abre la ventana, y
\texttt{py-slowdown ARCHIVO} la abre con ese archivo;
\texttt{py-slowdown -{}-help} muestra la ayuda, y \texttt{-{}-version} la
versión. Un archivo de audio también se puede soltar sobre la ventana, o
abrir con \menu{Archivo\then Abrir} (\key{Ctrl+O}) o \menu{Archivo\then
Abrir reciente}.

\section{La ventana}

\shot{window_es}{1}{La ventana, con una grabación a 0,75\,× y dos marcadores}

\index{ventana}\index{forma de onda}\index{mandos}\index{marcadores!panel}%
La forma de onda ocupa casi toda la ventana (figura~\ref{fig:window_es}).
Encima tiene una regla de tiempos; la línea roja es la posición, el punto
que se oye; lo ya tocado se ve más oscuro. Los marcadores son líneas de
trazos, con su nombre arriba. Debajo de la forma de onda están los mandos y
el tiempo, y debajo de ellos la velocidad, la altura y el volumen; a la
derecha, el panel de los
marcadores, que se puede cerrar, mover o dejar flotando
(\menu{Ver\then Marcadores} lo vuelve a mostrar). La barra de estado dice
qué se está leyendo, y luego el formato del archivo: su códec, su
frecuencia de muestreo, si es mono o estéreo, y su duración.

\section{Cómo se abre un archivo}

\index{formatos}\index{videos}%
El programa abre lo que lea FFmpeg: MP3, M4A y AAC, MP4, FLAC, WAV, AIFF,
Ogg, Opus, WMA y muchos más, y el sonido de los videos (MP4, MKV, WebM,
MOV). Un archivo con más de dos canales se mezcla en estéreo.

\index{archivos largos}\index{lectura}%
El archivo se lee en segundo plano. Mientras tanto, la barra de estado
muestra cuánto va, la forma de onda se llena desde la izquierda y lo ya
leído se puede tocar. Una grabación de dos horas tarda unos segundos. Un
archivo largo en MP3, FLAC, Ogg Vorbis u Opus se lee en cuatro partes a la
vez, cada una en un proceso propio: la forma de onda llena su primer cuarto
desde la izquierda, y luego el resto de una vez. El
sonido decodificado se guarda en un archivo temporal de la carpeta de
caché, no en la memoria (dos horas en estéreo ocupan allí unos 1,3\,GB), y
ese archivo desaparece al cerrar el archivo de audio, o al terminar el
programa.

\index{posicion recordada@posición recordada}%
Al cerrar un archivo (o al abrir otro, o al salir), el programa guarda el
punto donde iba, y la vez siguiente que se abre el archivo sigue desde
allí. Un archivo que se dejó en su final vuelve a empezar desde el
principio.

\section{Reproducción}

\shot{controls_es}{1}{Los mandos}

\index{mandos}\index{reproduccion@reproducción}\index{teclas}%
Los mandos (figura~\ref{fig:controls_es}) son, de izquierda a derecha:

\begin{description}
\item[Ir al principio] (\key{Inicio}).
\item[Atrás] la longitud del salto (\key{←}); con el botón apretado sigue
  retrocediendo, para recorrer el archivo.
\item[Tocar y pausar] (\key{Espacio}).
\item[Detener] y volver al principio (\key{S}).
\item[Adelante] la longitud del salto (\key{→}); con el botón apretado
  sigue avanzando.
\item[Ir al final] (\key{Fin}).
\item[Marcador] pone uno en la posición (\key{B}).
\item[Repetir] la selección, una y otra vez (\key{L}; sección~\ref{sec:repetir}).
\item[Cuenta previa] antes de tocar (\key{C}; sección~\ref{sec:cuenta}).
\end{description}

\noindent\index{longitud del salto}%
El salto es de 5 segundos, salvo que se elija otro en
\menu{Configuración\then Longitud del salto} (de 1 a 60 segundos).
\key{Mayús+←} y \key{Mayús+→} retroceden y avanzan un segundo, que es lo
que más hace falta al sacar un pasaje. Al final del archivo la reproducción se
detiene; \key{Espacio} vuelve entonces a tocarlo desde el principio.

\index{volumen}%
El volumen es la barra de la derecha. Solo baja el sonido del programa;
nunca lo sube por encima del del archivo.

\section{La velocidad}

\index{velocidad}%
La velocidad va de 0,25\,× (cuatro veces más lento) a 3\,× (tres veces más
rápido). Se cambia con la barra, escribiéndola en la casilla de al lado o
con las teclas \key{↓} y \key{↑} (o \key{−} y \key{+}), de 0,05 en 0,05;
\key{N}, o el botón \menu{1×}, vuelve a la velocidad normal. Se puede
cambiar mientras suena; el cambio se oye después de unos 0,15 segundos, el
sonido que ya estaba preparado a la velocidad anterior.

\index{calidad}\index{rubber band@Rubber Band}%
La altura no cambia: una melodía a 0,5\,× suena más despacio, pero en su
propio tono. La calidad es la del motor R3 de Rubber Band, que mantiene
natural la música, aun lejos de 1\,×; a velocidades muy lentas los ataques
de los sonidos (los golpes de una guitarra o de un tambor) se suavizan un
poco, como con cualquier método de esta clase. A
1\,× el archivo suena tal cual.

\index{velocidad!no disponible}%
Si la biblioteca Rubber Band no está instalada (lo que solo puede pasar
cuando el programa se instala desde el código), la ventana lo dice en la
barra de estado, y la velocidad se queda en 1\,×.

\section{La altura}
\label{sec:altura}

\index{altura}\index{semitonos|see{altura}}\index{transportar|see{altura}}%
La altura se puede subir o bajar, en semitonos, sin cambiar la velocidad:
para tocar con una grabación en otra tonalidad, para llevar una canción al
alcance de una voz, o para oír un pasaje rápido más lento y una octava más
arriba, donde sus notas se distinguen. La casilla \menu{Altura}, junto a
la velocidad, va de 12 semitonos abajo a 12 arriba (una octava a cada
lado), de uno en uno; \key{Ctrl+↑} y \key{Ctrl+↓} hacen lo mismo, y
\key{Mayús+N}, o el botón \menu{0}, vuelve a la altura de la grabación. La
casilla acepta también décimas, escritas en ella (0,5, por ejemplo): una
grabación algo desafinada, como suelen estarlo las cintas viejas, se puede
afinar con el instrumento.

\index{calidad}\index{rubber band@Rubber Band}%
La altura se puede cambiar mientras suena la grabación, sin chasquidos.
También la cambia Rubber Band, en la misma pasada que la velocidad y con la
misma calidad: medido, un tono queda a menos de un centésimo de semitono
de la altura pedida, también a las velocidades más lentas y más
rápidas. A la altura de la grabación y a 1\,× el archivo
suena tal cual. La altura, como la velocidad, no se guarda de una sesión a
otra: el programa empieza siempre en la de la grabación.

\section{La forma de onda}

\index{forma de onda}\index{posicion@posición}%
Un clic en la forma de onda lleva allí la posición, y arrastrar la mueve,
suene o no la grabación. La forma de onda va a la escala del punto más
fuerte del archivo, para que una grabación baja no se vea plana.

\shot{zoom_es}{0.85}{Una vista de cerca, alrededor de un marcador}

\index{acercar y alejar}\index{seguir la posicion@seguir la posición}%
\key{Ctrl+rueda} acerca o aleja la vista alrededor del puntero
(figura~\ref{fig:zoom_es}); \key{Ctrl++} y \key{Ctrl+−} hacen lo mismo
alrededor de la posición, y \key{Ctrl+0} vuelve a mostrar todo el archivo.
De cerca se dibujan las muestras mismas, hasta una por píxel. Cuando la
vista no muestra todo el archivo, aparece una barra de desplazamiento
debajo, y la rueda corre la vista; mientras suena, la vista pasa a la
página siguiente cuando la posición sale de ella, salvo que se desactive
\menu{Ver\then Seguir la posición}.

\section{Los marcadores}

\index{marcadores}%
Un marcador señala un punto del archivo, con un nombre si se quiere.
\key{B}, el botón \menu{Marcador} o \menu{Marcadores\then Añadir un
marcador} ponen uno en la posición; un clic derecho en la forma de onda
pone uno donde está el puntero (figura~\ref{fig:context_es}). Nunca se
ponen dos marcadores en el mismo punto: si ya hay uno, se elige ese.

\shot{context_es}{0.35}{El menú del clic derecho sobre un marcador}

\index{marcadores!nombre}%
El nombre se escribe en el panel, con un doble clic sobre él, o en el
diálogo que abre un doble clic sobre el marcador en la forma de onda, como
\menu{Cambiar el nombre del marcador} en el menú del clic derecho
(figura~\ref{fig:rename_es}). Un marcador sin nombre muestra su número.

\shot{rename_es}{0.35}{El nombre de un marcador}

\shot{panel_es}{0.45}{El panel de los marcadores}

\index{marcadores!panel}\index{marcadores!ir a}%
El panel (figura~\ref{fig:panel_es}) lista los marcadores en orden, con
sus tiempos. Un doble clic en un tiempo lleva allí la posición; un clic en
un marcador de la forma de onda, sobre su nombre, hace lo mismo.
\key{Ctrl+←} y \key{Ctrl+→} (o \key{RePág} y \key{AvPág}) van al marcador
anterior y al siguiente. Mientras suena, el anterior es el que está antes
del marcador que se acaba de pasar, de modo que un pasaje se puede volver a
oír desde su marcador con una sola tecla.

\index{marcadores!quitar}\index{marcadores!exportar}%
\menu{Quitar}, en el panel, o \menu{Quitar el marcador}, en el menú del
clic derecho, quita uno. \menu{Exportar}, en el panel, o
\menu{Marcadores\then Exportar los marcadores}, los guarda en un archivo de
texto, uno por línea: el tiempo, un tabulador y el nombre. Ese texto se
puede pegar en una partitura, en unas notas o en una hoja de cálculo.

\index{marcadores!donde se guardan@dónde se guardan}%
Los marcadores se guardan en cuanto cambian, por archivo, según su ruta
completa (sección~\ref{sec:archivos}). El archivo de audio nunca se
modifica. Un archivo que se mueve o se renombra es para el programa un
archivo nuevo, sin marcadores.

\section{Repetir una selección}
\label{sec:repetir}

\shot{repeat_es}{0.85}{Una selección que se repite, con la cuenta previa en su tercer golpe}

\index{seleccion@selección}\index{repetir}\index{bucle|see{repetir}}%
Un pasaje se puede repetir sin fin, para oírlo hasta entenderlo o para
tocar con él. Primero se selecciona: \key{Mayús+arrastrar} sobre la forma
de onda; o, con la posición en su principio, \key{I}, y en su final,
\key{O} (una sola de las dos selecciona de la posición al final del
archivo, o desde su principio); o, con el clic derecho,
\menu{Seleccionar de aquí al marcador siguiente}. La selección se dibuja
en verde, también sobre la regla, y la barra de estado dice dónde empieza y
dónde termina. Luego se activa \menu{Reproducción\then Repetir la
selección} (\key{L}), o el botón de las dos flechas: cuando la
reproducción llega al final de la selección, sigue desde su principio
(figura~\ref{fig:repeat_es}).

\index{repetir!sin pausa}%
Entre el final y el principio no hay pausa ni chasquido: la última
centésima de segundo de la selección se funde con lo que precede a su
principio, de modo que el pasaje suena como un solo flujo continuo. Esto
vale también a cualquier velocidad, porque el estiramiento se hace sobre
ese flujo, y no sobre cada repetición por separado. Un clic fuera de la
selección lleva allí la posición, como siempre; si es antes de la
selección, las repeticiones empiezan cuando la reproducción llega a ella.
\key{Esc}, o \menu{Quitar la selección}, la quita; \key{L} otra vez deja
de repetir y conserva la selección.

\section{La cuenta previa}
\label{sec:cuenta}

\shot{countin_es}{0.45}{Los ajustes de la cuenta previa}

\index{cuenta previa}\index{metronomo@metrónomo|see{cuenta previa}}\index{tempo}%
Antes de que empiece la grabación, el programa puede contar, como un
metrónomo antes de que entre una banda: unos golpes a un tempo, el primero
más agudo, y luego la grabación, en el pulso siguiente. Se activa con
\menu{Reproducción\then Contar antes de tocar} (\key{C}), o con el botón
\menu{Cuenta previa}; desde entonces cuenta cada vez que empieza a tocar
(con \key{Espacio} o con el botón), no cuando la posición cambia mientras
suena. Mientras cuenta, la forma de onda muestra el pulso en grande
(figura~\ref{fig:repeat_es}), y la posición no se mueve.

\index{cuenta previa!desfase}\index{desfase de la grabación}%
La grabación entra en el pulso que sigue al último golpe si su primer pulso
está justo en la posición. Si no lo está (un marcador puesto un poco antes
del primer pulso, por ejemplo), \menu{Desfase de la grabación}, en el
mismo diálogo, la mueve: un valor negativo la hace entrar antes, uno
positivo después, tantos segundos de la grabación (de 2 antes a 2 después,
en centésimas). Si el primer pulso está 0,2 segundos después de la
posición, $-0{,}2$ lo pone en el pulso, a cualquier velocidad. Solo puede
entrar antes hasta el último golpe, que nunca se corta: el diálogo avisa
cuando se pide más.

El desfase es más fácil de fijar sobre la forma de onda: con la posición
donde empezará a tocar, un clic derecho sobre el primer pulso de la
grabación, y \menu{Cuenta previa: el primer pulso está aquí}. El programa
fija el desfase de la posición a ese punto, activa la cuenta previa y dice
lo que hizo; una línea verde de trazos con un 1 muestra entonces, mientras
no suena, dónde espera la cuenta el primer pulso (figura~\ref{fig:zoom_es}),
para comprobarlo o para volver a fijarlo más de cerca. \menu{Cuenta previa:
el primer pulso en la posición} quita el desfase.

\index{marcar el tempo}\index{tempo!marcado}%
El tempo de una grabación se puede medir marcándolo: mientras suena, se
pulsa \menu{Marcar}, junto a \menu{Cuenta previa}, o la tecla \key{T},
con el pulso, unas cuantas veces. Desde el segundo golpe, la barra de
estado dice el tempo medido, que pasa a ser el de la cuenta previa; cuantos
más golpes, más estable, porque es el valor del medio de los intervalos
entre los ocho últimos, de modo que un golpe desigual no lo mueve. Una
pausa de más de dos segundos empieza una medida nueva. Lo que se marca es
el tempo que se oye: con \menu{Seguir la velocidad}, un pulso marcado a
60 BPM mientras la grabación suena a 0,5\,× se guarda como 120 BPM, el de
la grabación, y así la cuenta previa queda bien a cualquier velocidad. El
diálogo de la cuenta previa también tiene un botón \menu{Marcar}, junto al
tempo.

\index{bluetooth@Bluetooth}%
El primer golpe llega cuatro décimas de segundo después de Tocar, y la
cuenta empieza con un soplo muy tenue, apenas audible, que sigue entre los
golpes: muchos auriculares Bluetooth se silencian mientras reciben
silencio, y tardan un momento en despertar, de modo que, sin él, se
tragaban los golpes cortos de la cuenta.

La flecha junto al botón \menu{Cuenta previa}, \menu{Configuración\then
Cuenta previa} o \menu{Reproducción\then Ajustes de la cuenta previa}
abren sus ajustes (también un clic derecho sobre el botón), y la etiqueta
del botón dice cuáles son ahora. Fijan cuántos golpes (de 1 a 16) y el tempo, en pulsos por
minuto (de 20 a 300; figura~\ref{fig:countin_es}). Con \menu{Seguir la
velocidad}, el tempo cambia con la velocidad: una canción a 120 BPM tocada
a 0,5\,× se cuenta a 60, como va a sonar. El diálogo dice lo que se oirá a
la velocidad de ese momento. Estos ajustes se conservan.

\section{Configuración}

\shot{menu_settings_es}{0.35}{El menú Configuración}

\index{configuracion@configuración}\index{idioma de la interfaz}\index{longitud del salto}%
El menú \menu{Configuración} (figura~\ref{fig:menu_settings_es}) tiene los
ajustes de la cuenta previa (sección~\ref{sec:cuenta}), la longitud del
salto de \menu{Atrás} y \menu{Adelante} y el idioma de la
interfaz. El idioma sigue al del sistema, salvo que se elija uno; uno nuevo
se usa cuando el programa vuelve a abrirse. El programa recuerda además el
tamaño y el lugar de la ventana y de su panel, el volumen, si la vista
sigue la posición, la carpeta del último archivo abierto y los archivos
recientes.

\index{reabrir el último archivo}%
Con \menu{Configuración\then Reabrir el último archivo al empezar} marcado
(al principio no lo está), el programa vuelve a abrir, al empezar, el
último archivo que estuvo abierto, y sigue desde donde se dejó. Si se nombra
un archivo en la línea de órdenes, se abre ese; y si el último se movió o se
borró, el programa empieza sin archivo.

\section{Los archivos que guarda}
\label{sec:archivos}

\index{archivos de configuracion@archivos de configuración}\index{marcadores!donde se guardan@dónde se guardan}%
\begin{description}
\item[La configuración] en
  \texttt{\textasciitilde/.config/py-slowdown/py-slowdown.conf} en Linux; en
  el registro, bajo \texttt{HKEY\_CURRENT\_USER\textbackslash
  Software\textbackslash py-slowdown}, en Windows, y en
  \texttt{\textasciitilde/Library/Preferences/com.py-slowdown.py-slowdown.plist}
  en macOS. Si se borra, con el programa cerrado, vuelven los valores de
  siempre.
\item[Los marcadores] y el punto donde se dejó cada archivo, en
  \texttt{bookmarks.json}: en \texttt{\textasciitilde/.local/share/py-slowdown/}
  en Linux (en \texttt{\$XDG\_DATA\_HOME}, si está definida), en
  \texttt{\%APPDATA\%\textbackslash py-slowdown} en Windows y en
  \texttt{\textasciitilde/Library/Application Support/py-slowdown/} en macOS.
  Es un archivo JSON, que se puede copiar a otro computador o leer con
  otros programas.
\item[El sonido decodificado] del archivo abierto, en un archivo temporal
  sin nombre, que se va con él: en \texttt{\textasciitilde/.cache/py-slowdown/}
  en Linux, \texttt{\%LOCALAPPDATA\%\textbackslash py-slowdown\textbackslash
  cache} en Windows y \texttt{\textasciitilde/Library/Caches/py-slowdown/} en
  macOS.
\end{description}

\noindent Desinstalar el programa nunca los quita.

\section{Teclas}

\index{teclas}%
\begin{longtable}{@{}p{0.42\textwidth}p{0.52\textwidth}@{}}
\toprule
\textbf{Teclas} & \textbf{Qué hacen}\\
\midrule
\endhead
\key{Espacio} & tocar o pausar\\
\key{S} & detener y volver al principio\\
\key{Inicio}, \key{Fin} & ir al principio, al final\\
\key{←}, \key{→} & atrás, adelante (la longitud del salto)\\
\key{Mayús+←}, \key{Mayús+→} & un segundo atrás, adelante\\
\key{↓}, \key{↑} (o \key{−}, \key{+}) & más lento, más rápido (de 0,05\,× en 0,05\,×)\\
\key{N} & velocidad normal\\
\key{Ctrl+↑}, \key{Ctrl+↓} & un semitono más agudo, más grave\\
\key{Mayús+N} & la altura original\\
\key{B} (o \key{Ctrl+B}) & añadir un marcador en la posición\\
\key{Ctrl+←}, \key{Ctrl+→} (o \key{RePág}, \key{AvPág}) & marcador anterior, siguiente\\
\key{Ctrl++}, \key{Ctrl+−}, \key{Ctrl+0} & acercar, alejar, todo el archivo\\
\key{Ctrl+rueda} & acercar o alejar alrededor del puntero\\
\key{L} & repetir la selección\\
\key{Mayús+arrastrar} & seleccionar\\
\key{I}, \key{O} & principio, final de la selección en la posición\\
\key{Esc} & quitar la selección\\
\key{C} & contar antes de tocar\\
\key{T} & marcar el tempo, con el pulso\\
\key{Ctrl+O} & abrir un archivo\\
\key{Ctrl+W} & cerrar el archivo\\
\key{F1} & la lista de las teclas\\
\bottomrule
\end{longtable}

\section{Cuando algo falla}

\index{problemas}%
\begin{description}
\item[No suena nada] revise el volumen del programa y el del sistema, y
  que esté elegida la salida correcta en la configuración del sonido del
  sistema; el programa toca por la salida predeterminada.
\item[La velocidad o la altura no se pueden cambiar] la biblioteca Rubber Band no está
  instalada (la barra de estado lo dice): instálela como en la
  sección~\ref{sec:instalacion}.
\item[Un archivo no se abre] FFmpeg no pudo leerlo: puede estar dañado, o
  no ser audio. Un archivo dañado a la mitad se abre hasta donde se puede
  leer, y la barra de estado lo dice.
\item[Los marcadores de un archivo desaparecieron] el archivo se movió o se
  renombró; sus marcadores siguen en \texttt{bookmarks.json}, con su ruta
  anterior, y allí se pueden pasar a la nueva, con el programa cerrado.
\item[Un ajuste quedó mal] borre la configuración
  (sección~\ref{sec:archivos}).
\item[macOS o Windows no abren el programa] no está firmado por Apple ni
  por Microsoft: véase la sección~\ref{sec:instalacion}.
\end{description}

\noindent
Los problemas y los deseos se pueden escribir como incidencias en GitHub o
en Codeberg (las direcciones están en la primera página), o enviar a
\href{mailto:nicolasvaughan@proton.me}{\texttt{nicolasvaughan@proton.me}}.

\printindex

\end{document}

% ===============================================
%%% Local Variables:
%%% mode: latex
%%% lsp-ltex-plus-language: "es-CO"
%%% TeX-master: t
%%% End:
